% GENERATED FILE. Edit canonical lessons, then run npm run generate:books.
% canonical-source-hash: fnv1a64:43609c845e2282de
% canonical-lessons: KA-C292-nirjana, KA-C292-garigari, KA-C292-accukattada, KA-C292-astavyasta, KA-C292-anukula

\chapter{Deserted, Crisp, Tidy, Messy, Convenient}
\label{ch:ka-deserted-crisp-tidy-messy-convenient}
\hlchaptermodality{292}

\begin{chapteropening}
\textbf{By the end of this chapter:} \emph{I can say deserted, crisp, tidy, messy and convenient.}

Complete the last lesson of chapter 292: I can say deserted, crisp, tidy, messy and convenient.
\end{chapteropening}

\section[\texorpdfstring{\kn{ನಿರ್ಜನ} (nirjana)}{ನಿರ್ಜನ (nirjana)}]{\kn{ನಿರ್ಜನ} (nirjana) — deserted}

\label{lesson:KA-C292-nirjana}



\begin{quote}
Before the new one: say the Kannada for wonderful, then the Kannada for interesting.
\end{quote}

\subsection*{You'll want to know: \kn{ನಿರ್ಜನ}}
\textbf{\kn{ನಿರ್ಜನ}} — \emph{nirjana} — \textquotedblleft{}deserted\textquotedblright{}.

\begin{grammarlens}[title={What you've built}]
One more word for this chapter.
\end{grammarlens}

\subsection*{Your turn}
\begin{itemize}
\raggedright
  \item \emph{Say it:} \emph{nirjana}
  \item \emph{Say it:} \emph{nirjana}, once more
  \item \emph{Say it:} say \emph{nirjana}
  \item \emph{Recall:} say the Kannada for wonderful, then the Kannada for interesting, then say \emph{nirjana} again
\end{itemize}

\begin{tcolorbox}[breakable,colback=teal!4,colframe=teal!35!black,title={Before you move on}]
What does \kn{ನಿರ್ಜನ} mean? (\textquotedblleft{}Deserted\textquotedblright{}.) Say it once more.
\end{tcolorbox}
\section[\texorpdfstring{\kn{ಗರಿಗರಿ} (garigari)}{ಗರಿಗರಿ (garigari)}]{\kn{ಗರಿಗರಿ} (garigari) — crisp}

\label{lesson:KA-C292-garigari}



\begin{quote}
Before the new one: say the Kannada for interesting, then the Kannada for deserted.
\end{quote}

\subsection*{You'll want to know: \kn{ಗರಿಗರಿ}}
\textbf{\kn{ಗರಿಗರಿ}} — \emph{garigari} — \textquotedblleft{}crisp\textquotedblright{}.

\begin{grammarlens}[title={What you've built}]
One more word for this chapter.
\end{grammarlens}

\subsection*{Your turn}
\begin{itemize}
\raggedright
  \item \emph{Say it:} \emph{garigari}
  \item \emph{Say it:} \emph{garigari}, once more
  \item \emph{Say it:} say \emph{garigari}
  \item \emph{Recall:} say the Kannada for interesting, then the Kannada for deserted, then say \emph{garigari} again
\end{itemize}

\begin{tcolorbox}[breakable,colback=teal!4,colframe=teal!35!black,title={Before you move on}]
What does \kn{ಗರಿಗರಿ} mean? (\textquotedblleft{}Crisp\textquotedblright{}.) Say it once more.
\end{tcolorbox}
\section[\texorpdfstring{\kn{ಅಚ್ಚುಕಟ್ಟಾದ} (accukaṭṭāda)}{ಅಚ್ಚುಕಟ್ಟಾದ (accukaṭṭāda)}]{\kn{ಅಚ್ಚುಕಟ್ಟಾದ} (accukaṭṭāda) — tidy, well-organised}

\label{lesson:KA-C292-accukattada}



\begin{quote}
Before the new one: say the Kannada for deserted, then the Kannada for crisp.
\end{quote}

\subsection*{You'll want to know: \kn{ಅಚ್ಚುಕಟ್ಟಾದ}}
\textbf{\kn{ಅಚ್ಚುಕಟ್ಟಾದ}} — \emph{accukaṭṭāda} — \textquotedblleft{}tidy, well-organised\textquotedblright{}.

\begin{grammarlens}[title={What you've built}]
One more word for this chapter.
\end{grammarlens}

\subsection*{Your turn}
\begin{itemize}
\raggedright
  \item \emph{Say it:} \emph{accukaṭṭāda}
  \item \emph{Say it:} \emph{accukaṭṭāda}, once more
  \item \emph{Say it:} say \emph{accukaṭṭāda}
  \item \emph{Recall:} say the Kannada for deserted, then the Kannada for crisp, then say \emph{accukaṭṭāda} again
\end{itemize}

\begin{tcolorbox}[breakable,colback=teal!4,colframe=teal!35!black,title={Before you move on}]
What does \kn{ಅಚ್ಚುಕಟ್ಟಾದ} mean? (\textquotedblleft{}Tidy, well-organised\textquotedblright{}.) Say it once more.
\end{tcolorbox}
\section[\texorpdfstring{\kn{ಅಸ್ತವ್ಯಸ್ತ} (astavyasta)}{ಅಸ್ತವ್ಯಸ್ತ (astavyasta)}]{\kn{ಅಸ್ತವ್ಯಸ್ತ} (astavyasta) — messy, in disarray}

\label{lesson:KA-C292-astavyasta}



\begin{quote}
Before the new one: say the Kannada for crisp, then the Kannada for tidy.
\end{quote}

\subsection*{You'll want to know: \kn{ಅಸ್ತವ್ಯಸ್ತ}}
\textbf{\kn{ಅಸ್ತವ್ಯಸ್ತ}} — \emph{astavyasta} — \textquotedblleft{}messy, in disarray\textquotedblright{}.

\begin{grammarlens}[title={What you've built}]
One more word for this chapter.
\end{grammarlens}

\subsection*{Your turn}
\begin{itemize}
\raggedright
  \item \emph{Say it:} \emph{astavyasta}
  \item \emph{Say it:} \emph{astavyasta}, once more
  \item \emph{Say it:} say \emph{astavyasta}
  \item \emph{Recall:} say the Kannada for crisp, then the Kannada for tidy, then say \emph{astavyasta} again
\end{itemize}

\begin{tcolorbox}[breakable,colback=teal!4,colframe=teal!35!black,title={Before you move on}]
What does \kn{ಅಸ್ತವ್ಯಸ್ತ} mean? (\textquotedblleft{}Messy, in disarray\textquotedblright{}.) Say it once more.
\end{tcolorbox}
\section[\texorpdfstring{\kn{ಅನುಕೂಲ} (anukūla)}{ಅನುಕೂಲ (anukūla)}]{\kn{ಅನುಕೂಲ} (anukūla) — convenient; convenience}

\label{lesson:KA-C292-anukula}



\begin{quote}
Before the new one: say the Kannada for tidy, then the Kannada for messy.
\end{quote}

\subsection*{You'll want to know: \kn{ಅನುಕೂಲ}}
\textbf{\kn{ಅನುಕೂಲ}} — \emph{anukūla} — \textquotedblleft{}convenient; convenience\textquotedblright{}.

\begin{grammarlens}[title={What you've built}]
One more word for this chapter.
\end{grammarlens}

\subsection*{Your turn}
\begin{itemize}
\raggedright
  \item \emph{Say it:} \emph{anukūla}
  \item \emph{Say it:} \emph{anukūla}, once more
  \item \emph{Say it:} say \emph{anukūla}
  \item \emph{Recall:} say the Kannada for tidy, then the Kannada for messy, then say \emph{anukūla} again
\end{itemize}

\begin{tcolorbox}[breakable,colback=teal!4,colframe=teal!35!black,title={Before you move on}]
What does \kn{ಅನುಕೂಲ} mean? (\textquotedblleft{}Convenient; convenience\textquotedblright{}.) Say it once more.
\end{tcolorbox}
